\section{Proposed experimental methodology}

\subsection{Optical system and characterization}
Record camera model, lens, working distance, field of view, pixel size on the
board, spectral filtering, exposure, gain, illumination wavelength, irradiance
uniformity and board placement. Measure the coating's actual fluorescence response.
Use an interlocked and UV-opaque enclosure.

\subsection{Dataset and ground truth}
Catalogue each \emph{physical PCB}, lot, coating process, inspection date,
lighting setup and repeated capture in the dataset register. Obtain expert-labelled
required coverage, forbidden regions and adjudicated defect polygons/masks. Keep
test PCBs and, where feasible, production lots wholly separate from training.

\subsection{Comparison methods}
\begin{enumerate}
\item \textbf{Classical}: correction, registration, threshold/morphology,
connected components, reference-mask comparison.
\item \textbf{Standalone learned}: an established lightweight segmentation
model trained and evaluated without post hoc reference-specific tuning.
\item \textbf{Hybrid}: learned coating probability, geometric registration,
required/forbidden reference regions and transparent decision rules.
\end{enumerate}

\subsection{Ablations and robustness}
Predefine tests for illumination angle/intensity, exposure, crop resolution,
image normalization, registration accuracy, coating batches, PCB shifts,
defect size and inference engine. Every candidate is evaluated on the same
frozen specimens and annotation protocol.

\subsection{Performance measures}
Report pixel-level intersection-over-union and Dice score, defect-level precision
and recall, sensitivity by defect size, PCB-level false acceptance and false
rejection, uncertainty intervals, and latency distributions. Measure camera
acquisition, transfers, preprocessing, inference, postprocessing and I/O
separately, alongside total end-to-end cycle time.

For binary defect decisions, report
\begin{equation}
\mathrm{Recall} = \frac{\mathrm{TP}}{\mathrm{TP}+\mathrm{FN}}, \qquad
\mathrm{Precision} = \frac{\mathrm{TP}}{\mathrm{TP}+\mathrm{FP}}.
\end{equation}
Fix positive/negative definitions and the defect matching criterion before
calculating these metrics. The acceptance threshold must be chosen using
training/validation data, not the final test set.

\subsection{Reporting gates}
No accuracy, novelty, line-throughput or return-on-investment claim is considered
validated until independent experimental data supports it. Log failed
experiments as well as successful configurations.
